\chapter{Final Model Comparison and Conclusions}
\label{chap:final_comparison}

In this final section, we consolidate the results obtained across all experimental setups. We compare the generative Gaussian models, the discriminative Logistic Regression (LR) models, and the Support Vector Machines (SVM) developed in this laboratory. 

The comparison is based on the Minimum DCF for the primary application ($\pi_T = 0.1$), representing the optimal discriminatory potential of each architecture, as well as the Actual DCF, which reflects the model's native score calibration.

\section{Summary of Results}

Table \ref{tab:final_master_comparison} summarizes the optimal configurations and best performances achieved by each model type across our evaluations.

\begin{table}[H]
\centering
\caption{Comprehensive Model Comparison ($\pi_T = 0.1$)}
\label{tab:final_master_comparison}
\resizebox{\textwidth}{!}{
\begin{tabular}{lcccc}
\toprule
\textbf{Model Architecture} & \textbf{Hyperparameters} & \textbf{minDCF} & \textbf{actDCF} & \textbf{Calib. Gap} \\
\midrule
\multicolumn{5}{l}{\textbf{Generative Models}} \\
\midrule
Standard MVG & - & 0.2629 & 0.3051 & 0.0422 \\
Naive Bayes & - & 0.2570 & 0.3022 & 0.0452 \\
Tied MVG (Linear) & - & 0.3628 & 0.4061 & 0.0432 \\
\midrule
\multicolumn{5}{l}{\textbf{Logistic Regression (Discriminative)}} \\
\midrule
Standard LR (Linear) & $\lambda = 10^{-2}$ & 0.3611 & 0.4568 & 0.0957 \\
Quadratic LR & $\lambda = 3.16 \times 10^{-2}$ & 0.2436 & 0.4972 & 0.2536 \\
\midrule
\multicolumn{5}{l}{\textbf{Support Vector Machines (Discriminative Margin)}} \\
\midrule
Linear SVM (Centered) & $C = 10^{-1}, K=1.0$ & 0.3582 & 0.9573 & 0.5991 \\
Polynomial SVM ($d=2$) & $C = 3.16 \times 10^{-2}$ & 0.2455 & 0.6291 & 0.3836 \\
Polynomial SVM ($d=4$) & $C = 3.16 \times 10^{-3}$ & \textbf{0.1767} & 0.3975 & 0.2208 \\
RBF Kernel SVM & $C = 10^2, \gamma = 10^{-1}$ & 0.1845 & 0.4367 & 0.2522 \\
\bottomrule
\end{tabular}}
\end{table}

\section{Analysis and Comparative Discussion}

\subsection*{Which models achieve the best results?}

If we evaluate strictly on optimal theoretical ranking capability (Minimum DCF), the high-capacity kernel SVMs absolutely dominate the dataset. The \textbf{Polynomial SVM ($d=4$)} achieves the global minimum DCF of \textbf{0.1767}, closely followed by the \textbf{RBF Kernel SVM} at \textbf{0.1845}. 

However, as discussed in the previous chapters, this extreme performance is likely inflated by the models' massive capacities successfully fitting specific local clusters in the validation set, bordering on overfitting.

A much more robust comparison lies between the models constrained to quadratic boundaries: the \textbf{Quadratic Logistic Regression (0.2436)} and the \textbf{Polynomial SVM with $d=2$ (0.2455)}. These two models perform almost identically. This mathematical convergence is expected: both models compute non-linear decision boundaries in the original feature space based on quadratic interactions, differing only in their loss functions (Log-Loss vs. Hinge Loss).

Finally, all models constrained to linear boundaries (Tied MVG, Standard LR, and Linear SVM) consistently plateau at a much higher minDCF around $0.358 - 0.363$. 

\subsection*{What kind of separation rules or distribution assumptions characterize these models?}

The models evaluated up to this point fall into three distinct statistical philosophies:

\begin{enumerate}
    \item \textbf{Generative Models (MVG, Naive Bayes):} These models make explicit, rigid assumptions about the underlying probability density of the data (specifically, that features follow a Normal distribution). They separate classes by finding the boundary where the posterior probabilities intersect. While structurally inflexible if the data is non-Gaussian, they are highly robust to noise and naturally produce well-calibrated probabilistic scores.
    \item \textbf{Discriminative Probabilistic Models (Logistic Regression):} These models make no assumptions about the data distribution. Instead, they assume the posterior probability can be modeled directly using a logistic sigmoid function. They optimize a hyperplane by minimizing the cross-entropy loss, seeking a balance between separating the classes and outputting smooth, probabilistic confidence scores.
    \item \textbf{Discriminative Geometric Models (SVM):} Support Vector Machines make absolutely no assumptions about probability densities or posterior distributions. They are purely geometric algorithms that seek to maximize the physical spatial margin between the closest opposing points (the support vectors) while minimizing a Hinge Loss. 
\end{enumerate}

\subsection*{How do the results relate to the characteristics of the dataset?}

The empirical results perfectly validate the physical characteristics of the dataset observed in our exploratory analyses:
\begin{itemize}
    \item \textbf{Covariance Mismatch:} The universal failure of all Linear models ($> 0.35$ minDCF) compared to the Quadratic/Kernel models ($< 0.25$ minDCF) confirms that the "Authentic" and "Impostor" classes have vastly different covariance structures and geometric spreads. Non-linear boundaries are strictly mandatory for this application.
    \item \textbf{Non-Gaussian Skewness:} The Quadratic LR and Quadratic SVM both outperform the standard Generative MVG model. This confirms that while the data is somewhat Gaussian-like, features exhibiting strong skewness (e.g., Features 5 and 6) violate the strict assumptions of the MVG. The discriminative models bypass this density assumption, extracting superior geometric separation directly from the gradients.
\end{itemize}

\section{Final Considerations on Score Calibration}

While the Support Vector Machines offer the absolute lowest Minimum DCFs, they suffer from catastrophic native calibration. 

Unlike Generative models (which compute exact likelihoods) or Logistic Regression (which computes log-odds), the raw score of an SVM $s(x) = \sum \alpha_i z_i k(x_i, x) + b$ represents a physical geometric distance from the separating hyperplane. This distance has no inherent probabilistic meaning. Consequently, trying to apply the theoretical Bayes Log-Likelihood threshold ($t = -\log \frac{\pi_T}{1-\pi_T}$) directly to SVM distances results in massive actual DCFs (e.g., $0.9573$ for the Linear SVM). 

This concludes our current evaluation phase: Generative models offer ready-to-use, natively calibrated probabilities at the cost of peak accuracy. High-capacity discriminative models (like SVMs and Quadratic LR) offer state-of-the-art ranking capability, but their scores are raw and uncalibrated. Deploying these high-performance discriminative models in a real-world security environment will strictly require an additional score calibration step (such as Platt Scaling) to align their geometric outputs with valid Log-Likelihood Ratios.